% apx_courb-errata.tex -- extracted from chapters/apx_b_errata.tex (round5, author decision 1.1).
% Included with % apx_courb-errata.tex -- extracted from chapters/apx_b_errata.tex (round5, author decision 1.1).
% Included with % apx_courb-errata.tex -- extracted from chapters/apx_b_errata.tex (round5, author decision 1.1).
% Included with % apx_courb-errata.tex -- extracted from chapters/apx_b_errata.tex (round5, author decision 1.1).
% Included with \input{tables/apx_courb-errata}. Float placement, caption and every value
% are byte-identical to the version that was inline. The extraction was verified by building the
% pre-extraction tree with the same prose edits applied and comparing the RENDERED TEXT of both PDFs
% (pypdfium2 text layer): 260,848 characters each, identical. [correction] These headers previously
% claimed a 'rendered PDF checksum' comparison. No checksum comparison was ever run, and it would
% not have matched: the PDF bytes changed (1,322,758 -> 1,323,852) because this commit also carries
% a prose edit. Caught in review; the claim now names the check that was actually performed.
\begin{table}[htb]
\caption{Content errata of the CoUrb 2026 article and the corrections applied in
Chapter~\ref{ch:courb} (English translation of the published text).}
\label{tab:apx:courb-errata}
\centering
\small
\begin{tabular}{@{}p{0.42\textwidth}p{0.52\textwidth}@{}}
\toprule
\textbf{Defect in the published article} & \textbf{Correction in the chapter} \\
\midrule
Results and conclusion state wins over the baseline in ``16 of the 21
evaluated combinations''. &
Audited count: 15 of the 21 combinations, with one additional technical tie
(\emph{Outdoors} in Florida, where the baseline mean exceeds the best variant by
0.02 percentage points, within one standard deviation). \\
\addlinespace
Introduction, results, and conclusion state ``average gains of 20 to 24
percentage points'' on the category task. &
Audited per-state means: gains of 20.2 to 22.0 percentage points, considering the
better of the two spatial encoders in each combination (the disclosure the audit
requires). \\
\addlinespace
Introduction states ``wins in most scenarios''. &
``outperforms the baseline in most scenarios'' (verdict-verb rule; claim
strength unchanged). \\
\bottomrule
\end{tabular}
\end{table}
. Float placement, caption and every value
% are byte-identical to the version that was inline. The extraction was verified by building the
% pre-extraction tree with the same prose edits applied and comparing the RENDERED TEXT of both PDFs
% (pypdfium2 text layer): 260,848 characters each, identical. [correction] These headers previously
% claimed a 'rendered PDF checksum' comparison. No checksum comparison was ever run, and it would
% not have matched: the PDF bytes changed (1,322,758 -> 1,323,852) because this commit also carries
% a prose edit. Caught in review; the claim now names the check that was actually performed.
\begin{table}[htb]
\caption{Content errata of the CoUrb 2026 article and the corrections applied in
Chapter~\ref{ch:courb} (English translation of the published text).}
\label{tab:apx:courb-errata}
\centering
\small
\begin{tabular}{@{}p{0.42\textwidth}p{0.52\textwidth}@{}}
\toprule
\textbf{Defect in the published article} & \textbf{Correction in the chapter} \\
\midrule
Results and conclusion state wins over the baseline in ``16 of the 21
evaluated combinations''. &
Audited count: 15 of the 21 combinations, with one additional technical tie
(\emph{Outdoors} in Florida, where the baseline mean exceeds the best variant by
0.02 percentage points, within one standard deviation). \\
\addlinespace
Introduction, results, and conclusion state ``average gains of 20 to 24
percentage points'' on the category task. &
Audited per-state means: gains of 20.2 to 22.0 percentage points, considering the
better of the two spatial encoders in each combination (the disclosure the audit
requires). \\
\addlinespace
Introduction states ``wins in most scenarios''. &
``outperforms the baseline in most scenarios'' (verdict-verb rule; claim
strength unchanged). \\
\bottomrule
\end{tabular}
\end{table}
. Float placement, caption and every value
% are byte-identical to the version that was inline. The extraction was verified by building the
% pre-extraction tree with the same prose edits applied and comparing the RENDERED TEXT of both PDFs
% (pypdfium2 text layer): 260,848 characters each, identical. [correction] These headers previously
% claimed a 'rendered PDF checksum' comparison. No checksum comparison was ever run, and it would
% not have matched: the PDF bytes changed (1,322,758 -> 1,323,852) because this commit also carries
% a prose edit. Caught in review; the claim now names the check that was actually performed.
\begin{table}[htb]
\caption{Content errata of the CoUrb 2026 article and the corrections applied in
Chapter~\ref{ch:courb} (English translation of the published text).}
\label{tab:apx:courb-errata}
\centering
\small
\begin{tabular}{@{}p{0.42\textwidth}p{0.52\textwidth}@{}}
\toprule
\textbf{Defect in the published article} & \textbf{Correction in the chapter} \\
\midrule
Results and conclusion state wins over the baseline in ``16 of the 21
evaluated combinations''. &
Audited count: 15 of the 21 combinations, with one additional technical tie
(\emph{Outdoors} in Florida, where the baseline mean exceeds the best variant by
0.02 percentage points, within one standard deviation). \\
\addlinespace
Introduction, results, and conclusion state ``average gains of 20 to 24
percentage points'' on the category task. &
Audited per-state means: gains of 20.2 to 22.0 percentage points, considering the
better of the two spatial encoders in each combination (the disclosure the audit
requires). \\
\addlinespace
Introduction states ``wins in most scenarios''. &
``outperforms the baseline in most scenarios'' (verdict-verb rule; claim
strength unchanged). \\
\bottomrule
\end{tabular}
\end{table}
. Float placement, caption and every value
% are byte-identical to the version that was inline. The extraction was verified by building the
% pre-extraction tree with the same prose edits applied and comparing the RENDERED TEXT of both PDFs
% (pypdfium2 text layer): 260,848 characters each, identical. [correction] These headers previously
% claimed a 'rendered PDF checksum' comparison. No checksum comparison was ever run, and it would
% not have matched: the PDF bytes changed (1,322,758 -> 1,323,852) because this commit also carries
% a prose edit. Caught in review; the claim now names the check that was actually performed.
\begin{table}[htb]
\caption{Content errata of the CoUrb 2026 article and the corrections applied in
Chapter~\ref{ch:courb} (English translation of the published text).}
\label{tab:apx:courb-errata}
\centering
\small
\begin{tabular}{@{}p{0.42\textwidth}p{0.52\textwidth}@{}}
\toprule
\textbf{Defect in the published article} & \textbf{Correction in the chapter} \\
\midrule
Results and conclusion state wins over the baseline in ``16 of the 21
evaluated combinations''. &
Audited count: 15 of the 21 combinations, with one additional technical tie
(\emph{Outdoors} in Florida, where the baseline mean exceeds the best variant by
0.02 percentage points, within one standard deviation). \\
\addlinespace
Introduction, results, and conclusion state ``average gains of 20 to 24
percentage points'' on the category task. &
Audited per-state means: gains of 20.2 to 22.0 percentage points, considering the
better of the two spatial encoders in each combination (the disclosure the audit
requires). \\
\addlinespace
Introduction states ``wins in most scenarios''. &
``outperforms the baseline in most scenarios'' (verdict-verb rule; claim
strength unchanged). \\
\bottomrule
\end{tabular}
\end{table}
